%% ===========================================================================
%%  Sensor Identity Does Not Transfer
%%  Elsevier CAS double-column template (cas-dc), as used for the previous
%%  submission.  Compile on Overleaf with pdfLaTeX.
%%
%%  Requires: cas-dc.cls + cas-common.sty (bundled with Overleaf's Elsevier CAS
%%  template).  Figures live in ./figures, TikZ diagrams in ./diagrams.tex.
%%
%%  Every number is produced by a script in ../src and stored in ../results;
%%  the provenance of each table is given in a comment above it.
%% ===========================================================================
\documentclass[a4paper,fleqn]{cas-dc}

\usepackage[numbers]{natbib}
\usepackage{amsmath,amssymb}
\usepackage{booktabs}
\usepackage{graphicx}
\usepackage{multirow}
\usepackage{tikz}
\usetikzlibrary{arrows.meta,positioning,calc,fit}

\graphicspath{{figures/}}

\newcommand{\macf}{macro-F\textsubscript{1}}

% The CAS class prints an "ORCID(s):" line in the first-page footnote block even when
% no ORCID is supplied. Suppress the empty label; delete this if ORCIDs are added.
\makeatletter
\AtBeginDocument{\@ifundefined{printorcid}{}{\renewcommand{\printorcid}{}}}
\makeatother

\input{diagrams}

\begin{document}

\let\WriteBookmarks\relax
\def\floatpagepagefraction{1}
\def\textpagefraction{.001}

\shorttitle{Sensor identity does not transfer}
\shortauthors{M. Kbaier et al.}

\title[mode=title]{Sensor Identity Does Not Transfer: Revisiting Test-Time
Adaptation for Cross-Building Human Activity Recognition}

\author[1]{Mariem Kbaier}
\ead{kbaiermeryem@gmail.com}

\author[1]{Ons Abderrahim}

\author[1]{Manar Amayri}

\author[1]{Nizar Bouguila}
\cormark[1]
\ead{nizar.bouguila@concordia.ca}

\affiliation[1]{organization={Concordia Institute for Information Systems
  Engineering, Concordia University},
  city={Montreal}, state={QC}, country={Canada}}

\cortext[1]{Corresponding author.}

\begin{abstract}
Test-time adaptation (TTA) is an appealing answer to cross-building deployment in smart
buildings: sensors run continuously, and labelled data from a new site never arrives. We
evaluate three TTA methods developed for images (DIGA, PETAL and PeTTA) on human activity
recognition and occupancy estimation across buildings, and find that how much they help is
governed by the quality of the model they start from and not by the method. Occupancy
estimation supplies the cleanest test, because the two transfer directions differ only in
the quality of the source checkpoint: against a sound one no method improves and the
normalisation-based ones lose ground, while against a checkpoint worse than chance the
same methods gain $27$ to $41$ MCC points. Adaptation gain here measures source-model
deficiency, not adaptation quality.

Reaching that comparison required measuring the cross-building gap itself. Sliding windows
with unit stride cause adjacent samples to share nineteen of twenty timesteps, so a random
train/validation split places a near-duplicate of almost every validation sample into
training; under that split a source model reaches $0.998$ \macf, and grouped by recording
day it reaches $0.53$. The window definition makes almost no difference to this: the
inflation comes from the split. Features matter as much. The $106$ per-sensor features
that dominate the usual design score \emph{below} stratified-random guessing on the target
in both directions, because a given sensor index is a different physical device in each
building. With causal windows that an online system can compute, and thirteen
window-level summary statistics that name no sensor, a model transfers to the second
apartment within $0.6$ \macf{} points of a model trained on that apartment's own days, at
matched training-set size and on the same held-out days. Entering the first apartment the
same comparison costs $21.3$ points. One of the two directions therefore contains
essentially no domain shift, which the conventional source-minus-target figure of $8.8$
points conceals by charging the transfer for a difference between the buildings.

Against that baseline the three methods behave differently, and each result traces to a
specific published equation, not to our port. PETAL's coefficient weights the likelihood,
so the value selected by tuning on target labels removes the learning signal entirely.
PeTTA gates its update rate on a Mahalanobis distance summed over $64$ embedding
dimensions; across buildings that sum is large enough that the gate saturates and the
teacher stops updating. DIGA does help where there is headroom, but modestly and conditionally: $1.7$ \macf{}
points with hyperparameters retuned on target labels, $0.3$ under a label-free criterion,
and across its full published grid the sign of the effect follows the direction of
transfer rather than the setting. Our conclusion is not that TTA fails in smart buildings, but that it was being
asked to repair a problem largely created by evaluation and feature design.
\end{abstract}

\begin{keywords}
Test-time adaptation \sep Human activity recognition \sep Occupancy estimation
\sep Smart buildings \sep Evaluation leakage
\end{keywords}

\maketitle

% ============================================================================
\section{Introduction}
\label{sec:intro}

Buildings account for a large share of global energy demand, and occupancy-aware control
of HVAC and lighting is one of the more direct levers available for reducing it
\cite{manfren2021,amangeldy2025}. Human activity recognition (HAR) and occupancy
estimation (OE) supply the signal such control needs \cite{dridi2022,amayri2019}. Both are
learned from sensor traces, and both face the same obstacle when a model trained in one
building is installed in another: room layouts differ, sensors sit in different places,
and occupants keep different routines \cite{shahbazian2023,hong2017}. A model that works
in the building where it was trained often does not work in the next one.

Test-time adaptation addresses this without requiring labels from the new site. A frozen
source model updates itself online from the unlabelled stream it is already processing,
with no access to source data and no target annotation \cite{liang2024}. For buildings,
where sensors are always on and manual annotation is impractical, the fit is obvious. The
method family has matured quickly on images: DIGA \cite{diga} adapts normalisation
statistics and class prototypes without gradients, PETAL \cite{petal} treats lifelong
adaptation as probabilistic inference, and PeTTA \cite{petta} stabilises adaptation when
environments recur. Whether their assumptions survive the move to one-dimensional sensor
data is an open question, and the one we set out to answer.

Answering it requires first establishing that the quantity these methods are asked to
reduce is measured soundly, and we began there. Implementing the standard pipeline
reproduces the picture the literature reports: a source model validating near $0.99$
\macf{} and scoring $0.16$ on a second apartment, a drop that reads as severe domain
shift. That measurement does not survive scrutiny. Features are computed over sliding
windows with unit stride, so consecutive samples share $19$ of their $20$ timesteps, and a
random train/validation split places a $95\%$-identical copy of nearly every validation
sample into training. We isolate the effect by varying the window definition and the split
independently: the window moves source validation by under two points, the split moves it
by forty-six. Grouping by recording day, a one-line change, gives $0.53$ and leaves
the target number essentially untouched. The gap that motivates this literature is
therefore substantially smaller than reported, and the remainder of this paper measures
both domains under the same protocol.

Once the protocol is fixed, the picture changes in a way that is more interesting than
the original. Most of the remaining gap is attributable not to some irreducible property
of cross-building transfer but to the features themselves. The conventional design devotes
$106$ of its $135$ dimensions to per-sensor statistics, and those features are actively
harmful across buildings: used alone they score below chance in both transfer directions,
because \texttt{sensor\_id\_100} is a different physical device in each apartment. A model
built from thirteen layout-invariant window statistics transfers far better while
validating slightly worse in-domain.

With the features fixed, the question of what cross-building transfer actually costs can
be put properly. Source-validation minus target performance, the quantity usually quoted,
compares one building with another and so moves whenever the two differ in difficulty,
whether or not anything transfers badly. Holding the evaluation set fixed and varying only
where the training labels came from gives a different answer, and an asymmetric one. Into
the second apartment, a model trained on the first scores within $0.6$ points of a model
trained on the second apartment's own days, at matched training-set size: there is
essentially no cross-building penalty to speak of. Into the first apartment the same
comparison costs $21.3$ points. The price of changing building is not one number, and in
one of these two directions it is close to zero.

That asymmetry organises everything that follows. Where there is no penalty there is
nothing for adaptation to recover, and every method we test loses ground; where the
penalty is real, the one method that can act within a short stream gains. The three
methods otherwise divide cleanly. One helps a little. Two do not move at all, and in both
cases the reason is specific, mechanical, and traceable to the original paper rather than
to our implementation.

Our contributions are as follows.

\begin{itemize}
\item \textbf{A cross-building penalty measured against the right ceiling, and found to
  be asymmetric.} Holding the evaluation set fixed and matching training-set size, the
  penalty is $0.6$ points entering one apartment and $21.3$ entering the other
  (Section~\ref{sec:res-gap}). Per-sensor features are what create the apparent gap:
  used alone they score below stratified-random on the target in both directions.

\item \textbf{Adaptation gain measures source-model deficiency, not adaptation quality.} On cross-office
  occupancy estimation the two transfer directions differ only in the quality of the
  source checkpoint: against a sound one no method improves, against one worse than
  chance the same methods gain $27$ to $41$ MCC points. PETAL and PeTTA reproduce source
  predictions in both (Section~\ref{sec:res-oe}).

\item \textbf{Three mechanisms by which image-domain assumptions break on sensor data},
  each tied to a published equation: PETAL's $\bar\lambda$ scales the likelihood, PeTTA's
  divergence signal saturates, and DIGA's advertised global hyperparameters do not
  transfer (Section~\ref{sec:res-mech}).

\item \textbf{A protocol error that accounts for most of the reported cross-building gap.}
  We isolate it by varying the window definition and the validation split independently.
  Causality and segmentation move source \macf{} by less than two points; the split moves
  it by forty-six (Section~\ref{sec:res-leak}).

\item \textbf{Causal features that are both deployable and better.} Windows centred on the
  labelled sample read five minutes into the future, contradicting the online premise of
  TTA. Replacing them with causal windows improves target \macf{} by $7.3$ points. There
  is no trade-off to manage (Section~\ref{sec:res-leak}).

\item \textbf{A located ceiling.} Seven attempts to raise target performance beyond the
  deployable checkpoint are reported, six of which fail, including an unsupervised
  sensor-correspondence method given the advantage of offline access to the target
  stream (Sections~\ref{sec:res-ceiling} and~\ref{sec:res-corr}).
\end{itemize}

Section~\ref{sec:related} reviews domain adaptation and TTA for sensor data.
Section~\ref{sec:background} states the problem and the three methods.
Section~\ref{sec:method} describes our data handling, protocol and features.
Sections~\ref{sec:setup} and~\ref{sec:results} give the setup and results.
Sections~\ref{sec:discussion} and~\ref{sec:conclusion} discuss and conclude.

% ============================================================================
\section{Related work}
\label{sec:related}

\subsection{Domain adaptation in smart buildings}

Cross-environment degradation in building sensor models is well documented, and the usual
remedy is feature-level alignment using unlabelled target data gathered before deployment
\cite{dridi2022,abderrahim2025,amayri2019}. Open-set extensions handle activities present
in the target but absent from source training \cite{abderrahim2026,abderrahim_osda}. What
these approaches share is an offline phase: target data must be collected and processed
before the model goes live. A model arriving in a new building has none, annotating sensor
streams is expensive, and taking a building management system offline to retrain is rarely
acceptable. TTA exists to remove that phase.

\subsection{Test-time adaptation and its assumptions}

TTA adapts a pre-trained model during inference using only unlabelled target samples
\cite{liang2024}. It is stricter than source-free domain adaptation \cite{shot}, which
permits repeated offline passes over the target set: TTA runs online, in a single pass, on
samples as they arrive. TENT \cite{tent} minimises prediction entropy over batch
normalisation parameters and assumes a stationary, i.i.d.\ test stream, both of which
fail for building sensors, where distributions drift across the day and consecutive
readings are correlated. Under class imbalance, entropy minimisation tends to reinforce
the majority class, since confident low-entropy outputs are easiest to produce for
frequent classes. CoTTA \cite{cotta} extends TENT to continual settings with weight
averaging and stochastic restoration, but relies on augmentation-averaged pseudo-labels,
for which one-dimensional sensor data has no natural analogue.

The three methods we study were selected because their mechanisms are, at least in
principle, modality-agnostic. DIGA \cite{diga} mixes source and instance normalisation
statistics and adds a prototype classifier, computing no gradients at all. PETAL
\cite{petal} places a posterior over source weights and derives a student--teacher
objective with a source-posterior regulariser and Fisher-information-based restoration.
PeTTA \cite{petta} targets recurring environments (Fig.~\ref{fig:petta}), monitoring drift from the source model
and modulating its adaptation rate to prevent collapse. Section~\ref{sec:res-mech} shows
that each nonetheless embeds an assumption that fails here.

\subsection{TTA for one-dimensional data}

Outside vision, TTA is comparatively unexplored. Wang et al.~\cite{wang2023} propose an
optimisation-free method for wearable HAR using exponential-decay test-time
normalisation. Related work studies source-free adaptation for EEG, which shares temporal
correlation and channel heterogeneity with building sensors. None of this work addresses
the combination we face: severe class imbalance, differing sensor layouts across sites,
and recurring daily structure. Nor, to our knowledge, has any prior study examined how
evaluation protocol and feature design interact with measured adaptation gains ---
which, as we show, dominates the comparison.

% ============================================================================
\section{Background}
\label{sec:background}

\subsection{Problem formulation}

A domain is a pair $\mathcal{D}=(\mathcal{X},P(X))$ with feature space $\mathcal{X}$ and
input marginal $P(X)$. A model $f_\theta$ is trained on a source domain
$\mathcal{D}_S=(\mathcal{X}_S,P_S(X))$ from labelled data $\{(\mathbf{x}_n,y_n)\}_{n=1}^N$
and deployed on a target $\mathcal{D}_T$ with $P_S(X)\neq P_T(X)$ and no labels. At step
$t$ the model receives $\mathbf{x}_t\in\mathcal{X}_T$, predicts
$\hat y_t=f_{\theta_t}(\mathbf{x}_t)$, and updates $\theta_{t-1}\to\theta_t$ without
seeing a label:
\begin{equation}
\min_{\theta}\;
\mathbb{E}_{\mathbf{x}_t\sim P_T(X)}\big[\mathcal{L}_{\mathrm{TTA}}(f_\theta(\mathbf{x}_t))\big]
+\mathcal{R}(\theta,\theta_0),
\label{eq:tta}
\end{equation}
where $\theta_0$ are the source parameters and $\mathcal{R}$ penalises drift.

Two constraints on Eq.~\ref{eq:tta} are usually left implicit and matter a great deal
here. First, $f_{\theta_t}(\mathbf{x}_t)$ must be computable from information available at
time $t$; a feature that reads future samples cannot be deployed. Second, any quantity
used to define $\mathcal{X}$ (scalers, pruning masks, prototypes) must be estimable
without target labels, and ideally without target data. Section~\ref{sec:method} enforces
both.

\subsection{DIGA}

DIGA \cite{diga}, summarised in Fig.~\ref{fig:diga}, combines a Distribution Adaptation Module (DAM) and a Semantic Adaptation
Module (SAM), with no gradient computation. DAM mixes frozen source batch-normalisation
statistics with those of the current batch:
\begin{align}
\bar\mu^{T}_t &= \lambda_{\mathrm{BN}}\bar\mu^{S}+(1-\lambda_{\mathrm{BN}})\mu^{T}_t, \nonumber\\
(\bar\sigma^{T}_t)^2 &= \lambda_{\mathrm{BN}}(\bar\sigma^{S})^2+(1-\lambda_{\mathrm{BN}})(\sigma^{T}_t)^2 .
\label{eq:dam}
\end{align}
SAM estimates per-class prototypes $\mathbf{q}^c_t$ from predictions above confidence
$\mathcal{P}_0$, carries a historical copy by exponential moving average,
$\bar{\mathbf{q}}^c_t=\rho_P\bar{\mathbf{q}}^c_{t-1}+(1-\rho_P)\mathbf{q}^c_t$, and builds
\emph{two} non-parametric classifiers which are mixed by $\lambda_P$ before being fused
with the parametric output by $\lambda_F$:
\begin{align}
\tilde p &= \lambda_P\,p(c|\mathbf{x}_t,\bar{\mathbf{q}})+(1-\lambda_P)\,p(c|\mathbf{x}_t,\mathbf{q}), \\
p &= \lambda_F\,\tilde p+(1-\lambda_F)\,\hat p .
\label{eq:diga-fuse}
\end{align}
We note that DIGA's Eq.~8 writes $p(c|\mathbf{x},\mathbf{q})\propto\exp(-\langle
\mathbf{z},\mathbf{q}_c\rangle)$. Read literally as an inner product this ranks the
\emph{least} similar prototype highest; we therefore treat $\langle\cdot,\cdot\rangle$ as
a dissimilarity and report the effect of the alternatives in
Section~\ref{sec:res-mech}.


\begin{figure}[t]
\centering
\resizebox{0.95\columnwidth}{!}{\figDIGA}
\caption{DIGA applied to the one-dimensional backbone. DAM mixes frozen source
batch-normalisation statistics with the current batch's; SAM builds instance and
historical prototypes and fuses their predictions with the parametric classifier. No
gradient is computed anywhere in the loop.}
\label{fig:diga}
\end{figure}

\subsection{PETAL}

PETAL \cite{petal}, shown in Fig.~\ref{fig:petal}, approximates the source posterior $q(\theta)$ by SWAG-diagonal and
updates a student by ascending
\begin{equation}
\mathcal{L}=\log q(\theta)-\frac{\bar\lambda}{|\mathcal{B}|}\,
\mathrm{H}^{\mathrm{xe}}(y',y\,|\,\bar{\mathbf{x}}),
\qquad \theta_{t+1}=\theta_t+\eta\nabla_\theta\mathcal{L},
\label{eq:petal}
\end{equation}
with pseudo-labels $y'$ from an EMA teacher and a Fisher-information criterion restoring
unimportant parameters to $\theta_0$. The role of $\bar\lambda$ in Eq.~\ref{eq:petal} is
worth stating explicitly, because it is easy to misread: $\bar\lambda$ multiplies the
\emph{data} term, and the prior weight is fixed at unity. Setting $\bar\lambda=0$ does not
remove a regulariser; it removes the likelihood, leaving an update that does not depend on
the target stream at all.


\begin{figure*}[t]
\centering
\resizebox{0.80\textwidth}{!}{\figPETAL}
\caption{PETAL's student--teacher loop. The coefficient $\bar\lambda$ multiplies the
cross-entropy, not the source-posterior term; Section~\ref{sec:res-mech} shows what
follows when it is set to zero.}
\label{fig:petal}
\end{figure*}

\subsection{PeTTA}

PeTTA \cite{petta} targets recurring environments. It senses divergence in feature space
through a Mahalanobis distance between per-class running means and their source
counterparts,
\begin{equation}
\gamma^y_t=1-\exp\!\big(-(\hat\mu^y_t-\mu^y_0)^\top(\Sigma^y_0)^{-1}(\hat\mu^y_t-\mu^y_0)\big),
\label{eq:petta-gamma}
\end{equation}
and uses the batch average $\bar\gamma_t$ to set both the regularisation weight and the
teacher update rate,
\begin{equation}
\lambda_t=\bar\gamma_t\lambda_0,\qquad \alpha_t=(1-\bar\gamma_t)\alpha_0 .
\label{eq:petta-alpha}
\end{equation}
The quadratic form in Eq.~\ref{eq:petta-gamma} is a \emph{sum} over embedding dimensions.
Section~\ref{sec:res-mech} shows what that implies when the two domains are different
buildings rather than corruptions of the same image.


\begin{figure*}[t]
\centering
\resizebox{0.86\textwidth}{!}{\figPeTTA}
\caption{PeTTA for recurring streams. A single scalar $\bar\gamma_t$, derived from a
Mahalanobis distance \emph{summed} over the embedding dimension, sets both the
regularisation weight and the teacher update rate. As $\bar\gamma_t\to1$ the update rate
goes to zero.}
\label{fig:petta}
\end{figure*}

% ============================================================================
\section{Methodology}
\label{sec:method}

\subsection{Windowing}
\label{sec:method-window}

Binary sensor activations are aggregated over windows of $W=20$ consecutive timesteps at
$30$\,s resolution, giving a $10$-minute window per labelled sample. Two choices in this
step are usually made implicitly, and we make both explicit because they turn out to
matter.

\paragraph{Causality.} A window may be \emph{centred} on the labelled sample, spanning
$[i-W/2,\,i+W/2]$, or \emph{causal}, spanning $[i-W+1,\,i]$. The centred form is common
and reads five minutes into the future of every sample, which is incompatible with online
deployment. We use causal windows throughout and quantify the difference in
Section~\ref{sec:res-leak}.

\paragraph{Segmentation.} Windows must be truncated somewhere. Segmenting on
\emph{ground-truth activity changes} guarantees that no window straddles two activities,
but hands the extractor the label transitions it is meant to predict. We segment instead
on \emph{recording continuity}: a new segment begins at a day change or a gap in the
$30$\,s sampling. This uses no labels and is what a deployed system can determine. Windows
are truncated at segment starts instead of dropped, which is what an online system
genuinely has after a gap.

\subsection{Validation protocol}
\label{sec:method-split}

Because stride is one, adjacent samples share $W-1$ of $W$ timesteps. A random
train/validation split therefore places a near-duplicate of nearly every validation sample
into training, and the resulting score measures memorisation, not generalisation.
We use \textbf{GroupKFold over \texttt{year\_day}}: no recording day appears in both
halves, which removes window overlap and same-day correlation together. All source model
selection in this paper uses this protocol. Section~\ref{sec:res-leak} reports what each
alternative gives.

\subsection{Features}
\label{sec:method-feat}

The conventional $135$-dimensional vector comprises $106$ per-sensor statistics (mean and
standard deviation for each of the $53$ shared sensors), five window-level activity
statistics, $20$ co-activation indicators for the most frequent sensor pairs in the
source, and four cyclic time encodings. Of these, $126$ are tied to \emph{sensor
identity}: the index $100$ refers to a different physical device in each building, so a
rule learned about it is not merely uninformative in the target but wrong.

We therefore evaluate a $13$-dimensional \textbf{layout-invariant} set: nine window-level
statistics of the per-timestep active-sensor count (mean, maximum, minimum, median,
standard deviation, idle fraction, activation entropy, transition count, transition rate)
and four cyclic encodings of time of day and day of week. Nothing in this set refers to
which sensor fired. Appendix~A reports a larger invariant set including a sorted
activation profile and run-length structure; it does not improve on the thirteen.

Fig.~\ref{fig:pca} projects the engineered space onto its first two source principal
components; the target occupies a displaced, compressed region.

All transforms are fitted on the source training split only. No target labels and no
target statistics are used at any point in constructing the source model, so every
configuration remains a legal TTA starting point.


\begin{figure}[t]
\centering
\resizebox{0.82\columnwidth}{!}{\figMLP}
\caption{The classifier used throughout. Only the input width changes between feature
sets; depth, widths and regularisation are held fixed so that comparisons isolate the
features.}
\label{fig:mlp}
\end{figure}

\begin{figure*}[t]
\centering
\includegraphics[width=0.92\textwidth]{fig_pca.pdf}
\caption{First two principal components of the causal engineered feature space, fitted on
S1 and applied unchanged to S2. The target occupies a displaced and markedly compressed
region, which is the shift the adaptation methods are asked to absorb.}
\label{fig:pca}
\end{figure*}

\subsection{Causal prediction smoothing}

Activities persist for minutes while predictions are made per $30$\,s window. We apply a
median filter over the last $L$ predictions, never crossing a recording-segment boundary.
This uses only past outputs and so remains deployable online. We report $L=41$, selected
on source data.

\subsection{Stream protocol}

Target data is divided into eight contiguous segments. \emph{Non-recurring} streams visit
each once; \emph{recurring} streams make three passes in randomised order, repeated over
three independent orderings. Target labels are withheld throughout adaptation and used
only to score predictions afterwards.

\subsection{What a method is allowed to know}
\label{sec:method-tiers}

Methods in this literature differ in what they require of the target building before they
can produce a prediction, and the distinction is usually left implicit. We separate three
levels.

\begin{description}
\item[Tier 1, no target data.] Everything is estimated on the source. The model can be
  installed in a new building and run immediately. This is the setting test-time
  adaptation exists to serve.
\item[Tier 2, unlabelled target data, offline.] The full target distribution must be
  collected before deployment, though no labels are needed. This is source-free domain
  adaptation. It is a weaker requirement than labels but it reintroduces exactly the
  data-collection phase that motivates TTA.
\item[Tier 3, target labels.] Not deployable.
\end{description}

The distinction matters for what follows. The layout-invariant features used throughout
this paper are tier 1: every transform is fitted on the source split alone, and the model
needs nothing from the target building. Feature-shift pruning, by contrast, selects
features by comparing source and target distributions, and is therefore tier 2 ---
unsupervised, but not deployable without a prior observation period. The hyperparameter
selection discussed next is tier 3.

\subsection{Hyperparameter selection}
\label{sec:method-hparam}

\paragraph{Two evaluation protocols, and which number belongs to which.} Results that
compare \emph{configurations} use the full protocol: GroupKFold over \texttt{year\_day},
$3$ seeds $\times$ $5$ folds, causal median smoothing at $L=41$, scored on the whole target
stream. Source-only there is \macf{} $0.4466$ (S1$\to$S2). Results that compare
\emph{adaptation methods} run a single stream over a single source checkpoint, since a
streaming method cannot be averaged over folds without restarting its state; those use the
last GroupKFold fold, no smoothing, and source-only is $0.4944$ averaged over the same
three seeds. The two differ because smoothing adds about $1$ point, fold choice a little
more, and the single-split source model has a seed spread of $0.027$. Tables~\ref{tab:lambda}
and~\ref{tab:diga} predate the three-seed convention and quote the seed-$0$ value of that
same checkpoint, $0.4641$; we give both so no figure in this paper is unattributable to a
protocol. Comparisons are always within a protocol, never across.

We report results using each method's \emph{published} hyperparameters, and separately
under target-label tuning where the two differ materially. Selecting on target labels is
tier 3 by the distinction above and is not available in deployment. It is also not merely
optimistic: Section~\ref{sec:res-mech} shows it selecting a configuration at which one
method cannot adapt at all, and reporting the maximum of a distribution that is mostly
negative for another. Where a label-free criterion is needed we use mean prediction
entropy on the target stream, and report what it selects alongside the oracle value. So
that no conclusion rests on a single setting, Section~\ref{sec:res-hgrid} searches each
method's full published hyperparameter grid on the activity task in both directions, and
Section~\ref{sec:res-mech} does the same for DIGA on occupancy.

% ============================================================================
\section{Experimental setup}
\label{sec:setup}

\subsection{Datasets}

\paragraph{HAR.} The dataset of Tapia et al.~\cite{tapia2004} provides two single-occupant
apartments, S1 and S2, instrumented with binary proximity sensors on household objects and
recorded over roughly two weeks each ($16$ recording days here). S1 exposes $70$ sensor columns and S2 exposes $75$; feature engineering uses the
$53$ shared by both. Table~\ref{tab:har-data} gives per-class counts. We evaluate both
transfer directions.

Class names follow the source annotation (\texttt{toileting}, \texttt{meal\_prep},
\texttt{watching\_tv}, \texttt{music}, \texttt{cleaning}). We note this because the
time-of-day distributions confirm them: activity code $65$ occurs between 07:31 and 14:51
with a median of 11:34, which is consistent with \texttt{watching\_tv} and not with a
sleep-type label.

\paragraph{Occupancy.} A private dataset collected from two university offices (H355 and
H358) at the Grenoble Institute of Technology using CO\textsubscript{2} concentration
sensors \cite{dridi_cluster,amayri2019}. The original annotation records occupancy counts
(zero, one, two, three, more than three); we binarise to occupied / not occupied. The
representation is nine engineered features produced by the originating pipeline; we use
them as supplied and perform no further feature construction, so no target information
enters through any step we control. Splits are $1145$ / $200$ source train / test and
$822$ / $150$ target train / test, with roughly $78\%$ / $22\%$ class balance. Because
the data is private, the released code reproduces the HAR results only.

\begin{table*}[t]
\centering
\caption{HAR dataset statistics. Class names follow the source annotation.}
\label{tab:har-data}
\small
\begin{tabular}{lrrrrrr}
\toprule
Split & Total & toileting & meal\_prep & watching\_tv & music & cleaning \\
\midrule
S1 & 5031 & 1065 & 857 & 1598 & 1069 & 442 \\
S2 & 3412 & 1266 & 417 & 1291 & 325 & 113 \\
\bottomrule
\end{tabular}
\end{table*}

\subsection{Model}

A three-layer MLP with batch normalisation, ReLU and dropout (Fig.~\ref{fig:mlp}),
reducing the input to a $64$-dimensional embedding followed by a five-class output layer. The architecture is held
fixed across all feature sets so that comparisons isolate the features.

\subsection{Metrics and baselines}

We report accuracy, \macf, MCC \cite{chicco2020} and per-class F\textsubscript{1}. \macf{} is the primary
metric for HAR because it weights all activities equally, which matters when the target
minority classes are small.

Chance baselines are reported in every table and are not optional here: several published
figures for this transfer setting sit at or below them. On S2, constant-majority
prediction gives accuracy $0.378$ and \macf{} $0.110$, and stratified-random prediction
gives accuracy $0.306$ and \macf{} $\mathbf{0.200}$. On S1 the stratified-random \macf{} is
$0.200$. For binary occupancy, always predicting ``not occupied'' gives accuracy $0.773$,
\macf{} $0.436$ and MCC $0.000$.

Differences between methods are reported with $95\%$ paired bootstrap confidence
intervals over the target stream ($2000$ resamples), computed on the difference from
source-only.

% ============================================================================
\section{Results}
\label{sec:results}

\subsection{Decomposing the source-domain fit}
\label{sec:res-leak}

Establishing what the cross-building gap actually is requires measuring both domains
under one protocol. Table~\ref{tab:leak} and Fig.~\ref{fig:leakage} vary the window
definition and the validation split independently, holding the feature arithmetic fixed,
which separates the contribution of each.

% source: results/exp4_windowing.csv  (src/exp4_windowing.py), 5 seeds
\begin{table*}[t]
\centering
\caption{Source-domain validation \macf{} by windowing policy and validation protocol
($135$ features, z-score, mean $\pm$ s.d.\ over 5 seeds). Reading down any column shows
that the window definition is nearly irrelevant; reading across any row shows that the
split is not.}
\label{tab:leak}
\small
\begin{tabular}{llccc}
\toprule
Causality & Segmentation & Random 80/20 & Blocked & Grouped by day \\
\midrule
centred & activity labels & $0.9984\pm0.0009$ & $0.4357\pm0.0345$ & $0.5287\pm0.0067$ \\
centred & recording       & $0.9979\pm0.0006$ & $0.4571\pm0.0451$ & $0.5258\pm0.0112$ \\
causal  & activity labels & $0.9975\pm0.0021$ & $0.4414\pm0.0131$ & $0.5261\pm0.0084$ \\
causal  & recording       & $0.9968\pm0.0022$ & $0.4608\pm0.0209$ & $0.5110\pm0.0115$ \\
\bottomrule
\end{tabular}
\end{table*}


\begin{figure*}[t]
\centering
\includegraphics[width=0.92\textwidth]{fig_leakage.pdf}
\caption{Source validation \macf{} under every combination of windowing policy and
validation split (mean $\pm$ s.d.\ over 5 seeds). The hatched bars are the random split.
Its height is independent of how the window is defined, which locates the inflation in the
split rather than in the features.}
\label{fig:leakage}
\end{figure*}

A fully causal pipeline with no label-derived segmentation still validates at $0.9968$
under a random split. The inflation is therefore attributable to window overlap under
random splitting, not to the window definition. Grouping by recording day gives
$0.51$--$0.53$, and we treat that as the honest source-domain figure.

Causality nonetheless matters, in the opposite direction to what one might expect. Under
a fixed training protocol, making features causal \emph{improves} target transfer
(Table~\ref{tab:window-target}): \macf{} rises from $0.1626$ to $0.2356$, with accuracy
rising from $0.375$ to $0.444$. The deployable pipeline is also the better-transferring
one, so nothing is traded away by adopting it.

% source: results/exp4_windowing.csv
\begin{table*}[t]
\centering
\caption{Target transfer (S1$\to$S2) by windowing policy, $135$ features, blocked-trained,
mean over 5 seeds. Stratified-random \macf{} is $0.200$.}
\label{tab:window-target}
\small
\begin{tabular}{llccc}
\toprule
Causality & Segmentation & \macf & MCC & Accuracy \\
\midrule
centred & activity labels & $0.1626\pm0.0125$ & $0.0945$ & $0.3753$ \\
centred & recording       & $0.1953\pm0.0274$ & $0.1327$ & $0.3924$ \\
causal  & activity labels & $0.2344\pm0.0232$ & $0.2036$ & $0.4533$ \\
causal  & recording       & $\mathbf{0.2356}\pm0.0124$ & $0.1805$ & $0.4440$ \\
\bottomrule
\end{tabular}
\end{table*}

\subsection{Feature design}
\label{sec:res-feat}

Table~\ref{tab:features} sweeps feature subsets and source-fitted transforms on causal,
recording-segmented features under GroupKFold selection.

% source: results/exp5_features_causal_recording.csv (src/exp5_features_causal.py)
\begin{table*}[t]
\centering
\caption{S1$\to$S2 transfer by feature subset and source-fitted transform (causal
features, GroupKFold over 16 days, 5 seeds). Stratified-random \macf{} is $0.200$.}
\label{tab:features}
\small
\begin{tabular}{llcccc}
\toprule
Feature set & Transform & $d$ & Source-val \macf & Target \macf & Target MCC \\
\midrule
All (conventional)    & z-score     & 135 & $0.5083$ & $0.2436$ & $0.1682$ \\
All                   & rank-gauss  & 135 & $0.5890$ & $0.3775$ & $0.3144$ \\
KS-pruned 20\%        & z-score     & 108 & $0.5164$ & $0.3130$ & $0.2090$ \\
KS-pruned 20\%        & rank-gauss  & 108 & $\mathbf{0.6232}$ & $0.3004$ & $0.2138$ \\
No per-sensor         & z-score     & 29  & $0.5121$ & $0.3561$ & $0.2633$ \\
Global + temporal     & z-score     & 9   & $0.5341$ & $\mathbf{0.4133}$ & $\mathbf{0.3490}$ \\
Global only           & z-score     & 5   & $0.2914$ & $0.2129$ & $0.0970$ \\
Per-sensor only       & z-score     & 106 & $0.3292$ & $0.1484$ & $0.0436$ \\
\bottomrule
\end{tabular}
\end{table*}

Two observations. Per-sensor features alone reach $0.1484$ \macf, which is $5.1$ points
\emph{below} stratified-random guessing. They are not weakly informative; they encode
rules that are wrong in the target building. And the ordering by source validation is not
the ordering by transfer: the best-validating row (KS-pruned with rank-gauss, $0.6232$)
transfers at $0.3004$, while the best-transferring row ($0.4133$) validates $8.9$ points
lower. Across the eight rows the two columns are still positively related
($\rho=+0.69$), so this is a reordering near the top and not a reversal; what it costs a
practitioner is quantified in Section~\ref{sec:res-gap}.

Table~\ref{tab:both-dir} and Fig.~\ref{fig:features} repeat the comparison in both
directions.

% source: results/exp11_reverse.csv (src/exp11_reverse.py), 3 seeds x 5 folds
\begin{table*}[t]
\centering
\caption{Forward and reverse transfer under an identical protocol (causal features,
GroupKFold over the source's days, causal smoothing $L=41$). ``Gap'' is source-validation
minus target \macf, in points.}
\label{tab:both-dir}
\small
\begin{tabular}{llcccccc}
\toprule
Direction & Feature set & $d$ & Source-val & Target \macf & MCC & Gap & vs.\ chance \\
\midrule
S1$\to$S2 & Layout-invariant & 13  & $0.5341$ & $\mathbf{0.4466}$ & $0.3973$ & $8.8$  & $+24.7$ \\
S1$\to$S2 & All (conventional) & 135 & $0.5133$ & $0.2118$ & $0.0976$ & $30.2$ & $+1.2$ \\
S1$\to$S2 & Per-sensor only  & 106 & $0.3248$ & $0.1350$ & $0.0123$ & $19.0$ & $\mathbf{-6.5}$ \\
S1$\to$S2 & No per-sensor    & 29  & $0.5119$ & $0.3630$ & $0.2734$ & $14.9$ & $+16.3$ \\
\midrule
S2$\to$S1 & Layout-invariant & 13  & $0.4538$ & $\mathbf{0.3068}$ & $0.2667$ & $14.7$ & $+10.7$ \\
S2$\to$S1 & All (conventional) & 135 & $\mathbf{0.5381}$ & $0.1583$ & $0.0332$ & $38.0$ & $\mathbf{-4.2}$ \\
S2$\to$S1 & Per-sensor only  & 106 & $0.3668$ & $0.1702$ & $0.0553$ & $19.7$ & $\mathbf{-3.0}$ \\
S2$\to$S1 & No per-sensor    & 29  & $0.4483$ & $0.2271$ & $0.1043$ & $22.1$ & $+2.7$ \\
\bottomrule
\end{tabular}
\end{table*}


\begin{figure*}[t]
\centering
\includegraphics[width=0.92\textwidth]{fig_features.pdf}
\caption{Target \macf{} by feature set in both transfer directions, with the
stratified-random baseline marked (dashed). Bars falling short of it are shown in red.
Per-sensor features are below chance in both directions; the conventional 135-d vector is
below chance in one of them.}
\label{fig:features}
\end{figure*}

The below-chance result replicates: per-sensor features fall short of stratified-random
guessing in both directions. The reordering is sharper in reverse, where the conventional
$135$-d set has the best source-validation score of any row ($0.5381$) and the worst
transfer ($0.1583$), itself $4.2$ points below chance. Choosing between these four feature
sets by source validation alone would select that row.

The gap is direction-dependent: $8.8$ points forward and $14.7$ reverse. The asymmetry has
a visible cause rather than being noise. S2 is the weaker source: $3412$ samples
against $5031$, with the hard classes much thinner ($113$ versus $442$ for
\texttt{cleaning}) --- and training on the smaller, more imbalanced apartment lowers both
sides of the comparison. We therefore state the result as a gap of $9$ to $15$ points
instead of quoting the more favourable direction.

\subsection{What does changing building actually cost?}
\label{sec:res-gap}

Table~\ref{tab:both-dir} reports source-validation minus target \macf{} and calls it a
gap, as is conventional. That quantity compares performance on one building with
performance on another, so it moves whenever the two differ in difficulty, whether or not
anything transfers badly. S1 and S2 do differ: $5031$ samples against $3412$, and $442$
\texttt{cleaning} windows against $113$. Part of the $8.8$ and $14.7$ points is the
buildings, not the transfer.

The quantity that isolates domain shift holds the evaluation set fixed and varies only
where the training labels came from. For each apartment in turn we train one model on
that apartment's own days and another on the other apartment entirely, and score both on
the same held-out days under GroupKFold, with identical features, folds and smoothing.
Training-set sizes are matched by subsampling, because S2 has two thirds of S1's data and
an unmatched comparison would charge domain shift for a difference in sample count.

% source: src/exp21_domain_gap.py -> results/exp21_matched.csv
\begin{table}[t]
\centering
\caption{Cross-building penalty on a fixed evaluation set at matched training-set size,
$3$ seeds $\times$ $5$ folds. \emph{Own days} trains on the target apartment; \emph{other
building} trains on the source apartment. The difference is what a deployment pays for
having no labels in the new building.}
\label{tab:gap}
\small
\begin{tabular}{@{}lccc@{}}
\toprule
Evaluated on & Own days & Other building & Penalty \\
\midrule
S2 held-out days & $0.3497$ & $0.3434$ & $\mathbf{0.6}$ pts \\
S1 held-out days & $0.4897$ & $0.2767$ & $21.3$ pts \\
\bottomrule
\end{tabular}
\end{table}

The penalty is not one number. Entering S2 it is $0.6$ points, which is well inside the
seed spread of either model and is not distinguishable from zero. A model that has never
seen S2 performs as well on S2 as a model trained on S2's own days. Entering S1 the same
comparison costs $21.3$ points.

Two things follow, and both matter for how the rest of this paper should be read.

The first is that the conventional gap figure overstates the forward direction badly. The
$8.8$ points of Table~\ref{tab:both-dir} are almost entirely the difference between the
two apartments as prediction problems, not the cost of moving between them. Reported
cross-building drops that take this form should be read the same way.

The second is that one of our two transfer directions contains no domain shift for a
method to correct. This is not a limitation of the experiment but the most useful thing in
it: it provides a control. A test-time adaptation method run on S1$\to$S2 is being asked
to improve a model that is already performing at its target ceiling, and any movement it
produces can only be downward. Run on S2$\to$S1 the same method has $21$ points of genuine
headroom. Section~\ref{sec:res-hgrid} shows the published methods behaving exactly along
that line, and Section~\ref{sec:res-oe} reproduces the same split on a second task.

\paragraph{Does in-domain validation predict transfer?} The asymmetry above raises the
related question of whether source-domain fit is a usable guide when choosing what to
deploy. Pooling every configuration in this study for which we have a matched
source-validation and target score ($70$ in total, across feature sets, transforms,
window lengths, classifier families and temporal models), the rank correlation is
$+0.63$ ($p<0.001$). In-domain validation is therefore positively informative overall,
and we make no claim to the contrary. Among configurations that are actually competitive,
however, it degrades: restricting to those validating at $0.48$ or above leaves
$\rho=+0.27$ ($p=0.11$), and the single best-validating configuration is not the best
transferring one in either direction of Table~\ref{tab:both-dir}. The honest summary is
that in-domain validation separates working configurations from broken ones and carries
little information beyond that, which is the regime a practitioner choosing between
reasonable alternatives is in.

\subsection{Test-time adaptation on the deployable checkpoint}
\label{sec:res-tta}

Table~\ref{tab:tta} compares the methods on the deployable checkpoint.

% source: results/exp6_tta_global+temporal.csv (src/exp6_tta_table.py), 3 seeds
\begin{table*}[t]
\centering
\caption{TTA on the layout-invariant causal checkpoint (3 seeds; recurring streams
additionally averaged over 3 randomised segment orderings). Stratified-random \macf{} is
$0.200$.}
\label{tab:tta}
\small
\begin{tabular}{llccc}
\toprule
Method & Stream & Accuracy & \macf & MCC \\
\midrule
Source-only & non-recurring & $0.6042\pm0.013$ & $0.4467\pm0.029$ & $0.4515\pm0.033$ \\
AdaBN       & non-recurring & $0.3782\pm0.011$ & $0.2977\pm0.005$ & $0.1517\pm0.013$ \\
TENT        & non-recurring & $0.3784\pm0.011$ & $0.2976\pm0.005$ & $0.1515\pm0.013$ \\
DIGA        & non-recurring & $0.6045\pm0.013$ & $\mathbf{0.4639}\pm0.001$ & $\mathbf{0.4739}\pm0.009$ \\
PETAL       & non-recurring & $0.6033\pm0.008$ & $0.4450\pm0.023$ & $0.4522\pm0.026$ \\
\midrule
Source-only & recurring & $0.6042\pm0.011$ & $0.4467\pm0.025$ & $0.4515\pm0.029$ \\
DIGA        & recurring & $\mathbf{0.6073}\pm0.010$ & $0.4625\pm0.013$ & $0.4700\pm0.022$ \\
PETAL       & recurring & $0.6038\pm0.010$ & $0.4431\pm0.020$ & $0.4504\pm0.026$ \\
PeTTA       & recurring & $0.6045\pm0.011$ & $0.4470\pm0.025$ & $0.4517\pm0.029$ \\
\bottomrule
\end{tabular}
\end{table*}

% source: results/exp6 bootstrap block
\begin{table*}[t]
\centering
\caption{$95\%$ paired bootstrap confidence intervals on the difference from source-only
($2000$ resamples, same stream).}
\label{tab:boot}
\small
\begin{tabular}{llcc}
\toprule
Method & Stream & $\Delta$\macf & $\Delta$MCC \\
\midrule
DIGA  & non-recurring & $[+0.0032,\,+0.0225]$ & $[+0.0277,\,+0.0559]$ \\
DIGA  & recurring     & $[+0.0146,\,+0.0311]$ & $[+0.0387,\,+0.0633]$ \\
PETAL & non-recurring & $[-0.0076,\,+0.0027]$ & $[-0.0013,\,+0.0130]$ \\
PETAL & recurring     & $[-0.0087,\,+0.0015]$ & $[-0.0067,\,+0.0067]$ \\
PeTTA & recurring     & $[-0.0007,\,+0.0006]$ & $[-0.0009,\,+0.0011]$ \\
AdaBN & non-recurring & $[-0.1754,\,-0.1379]$ & $[-0.3046,\,-0.2500]$ \\
TENT  & non-recurring & $[-0.1756,\,-0.1383]$ & $[-0.3050,\,-0.2503]$ \\
\bottomrule
\end{tabular}
\end{table*}

In Table~\ref{tab:boot}, DIGA is the only method whose interval excludes zero on the
positive side, and the effect is small: $+1.7$ \macf{} and $+2.2$ MCC points. It is also
contingent on how two under-specified steps of the method are implemented, which
Section~\ref{sec:res-digaimpl} quantifies. For comparison, the feature redesign is
worth $+17.0$ \macf{} points before any adaptation runs. PETAL and PeTTA are statistically
indistinguishable from the unadapted model. AdaBN and TENT are substantially
\emph{worse}, by around $17$ points of \macf{}, which is itself informative: adapting
normalisation statistics helps when the source model is poor and hurts once the features
are already distribution-robust.

Per-class F\textsubscript{1} (Table~\ref{tab:perclass}) shows four of five classes working.
DIGA's prototype correction does what it is designed to do, lifting \texttt{meal\_prep}
by eight points at a cost of seven on \texttt{toileting}. \texttt{cleaning}, with $113$
target samples, remains near zero for every method.

% source: results/exp6_tta_global+temporal.csv per-class columns
\begin{table*}[t]
\centering
\caption{Per-class F\textsubscript{1} on S2 (mean over seeds and orderings).}
\label{tab:perclass}
\small
\begin{tabular}{lccccc}
\toprule
Method & toileting & meal\_prep & watching\_tv & music & cleaning \\
\midrule
Source-only & $0.551$ & $0.493$ & $0.746$ & $0.435$ & $0.009$ \\
DIGA        & $0.478$ & $\mathbf{0.572}$ & $\mathbf{0.766}$ & $\mathbf{0.468}$ & $0.035$ \\
PETAL       & $0.550$ & $0.474$ & $0.746$ & $0.440$ & $0.014$ \\
PeTTA       & $0.552$ & $0.494$ & $0.746$ & $0.434$ & $0.009$ \\
AdaBN       & $0.365$ & $0.310$ & $0.482$ & $0.275$ & $0.056$ \\
\bottomrule
\end{tabular}
\end{table*}

\subsection{What does not move the result}
\label{sec:res-ceiling}

Target performance is bounded by source performance: the deployable checkpoint validates
at $0.534$ in-domain and transfers at $0.447$, a gap of under nine points. Raising the
target therefore requires raising the source with features that still cross buildings. We
report seven attempts, all negative, because together they locate the ceiling.

% source: results/exp10_improve.csv, exp15_classifier.csv, exp16_temporal.csv,
%         exp14_correspondence.csv
\begin{table*}[t]
\centering
\caption{Attempts to raise target \macf{} beyond the $13$-d causal checkpoint
(S1$\to$S2, same protocol throughout). Only causal smoothing and a source-estimated
transition model help, and both marginally.}
\label{tab:ceiling}
\small
\begin{tabular}{llcc}
\toprule
Attempt & What was varied & Source-val & Target \macf \\
\midrule
\textbf{baseline} & $13$-d, MLP, $W{=}20$, median $L{=}41$ & $0.5341$ & $\mathbf{0.4466}$ \\
\midrule
richer invariant features & $+$ sorted profile, runs, multi-scale ($60$-d) & $-$ & $0.4188$ \\
class weighting           & inverse-frequency weights          & $-$ & $0.3852$ \\
network width             & $64$--$32$ and $128$--$64$         & $-$ & $0.4297$ \\
classifier family         & boosting, forests, linear          & $0.5230$ & $0.4190$ \\
window length             & $W\in\{10,40,60\}$                 & $0.5611$ & $0.4330$ \\
temporal context          & stacked $T{=}11$, causal TCN $T{=}21$ & $0.5414$ & $0.4111$ \\
sensor correspondence     & unsupervised assignment (\S\ref{sec:res-corr}) & $-$ & $0.2227$ \\
\midrule
transition model          & source-estimated HMM, causal filtering & $0.5341$ & $0.4503$ \\
\bottomrule
\end{tabular}
\end{table*}

Two things are worth drawing out of Table~\ref{tab:ceiling}.

First, the only attempt that helps in \emph{both} transfer directions is the transition
model, and it helps by less than a point ($0.4466\to0.4503$ forward, $0.3068\to0.3160$
reverse). Its transition matrix is estimated on the source apartment and encodes nothing
about sensors, only how often one activity follows another. That it transfers at all,
while every attempt to extract more from the sensor stream does not, is consistent with
the rest of this paper: what crosses buildings is behavioural regularity, not layout.

Second, and more importantly, several rows raise the source-validation score while
lowering the target. Window length does ($0.5611$ at $W{=}60$ against $0.5341$ at
$W{=}20$), and so does temporal context ($0.5414$ for the TCN). Section~\ref{sec:res-feat}
observed the same thing across feature sets. The pattern is general, and
Section~\ref{sec:discussion} returns to it.

\subsection{Can the sensors be matched across buildings?}
\label{sec:res-corr}

If per-sensor features fail only because index $100$ denotes a different device in each
apartment, the obvious repair is to estimate the correspondence. We tested this. Each
sensor was described by a behavioural signature computed from unlabelled data: firing
rate, burst-length statistics, a $24$-bin diurnal profile, and an index-free summary of
its co-activation neighbourhood. The two sensor sets were then matched by balanced
assignment on signature distance.

To separate a real correspondence from overfitting, permutations were estimated on one
half of the target stream and scored on the other.

% source: results/exp14_correspondence.csv and the half-split control
\begin{table}[t]
\centering
\caption{Per-sensor features under four permutations, target \macf. A correspondence that
is real should transfer between halves of the target stream.}
\label{tab:corr}
\small
\begin{tabular}{lcc}
\toprule
Permutation & Fitted half & Held-out half \\
\midrule
identity (no matching)      & $0.0976$ & $0.1744$ \\
learned signatures          & $0.0749$ & $0.2348$ \\
greedy search on fitted half& $\mathbf{0.4174}$ & $0.1738$ \\
random, 20 draws            & --- & $0.1655\pm0.050$ \\
\bottomrule
\end{tabular}
\end{table}

A greedy search reaches $0.417$ on the half it is fitted to and $0.174$ on held-out data,
indistinguishable from identity and from random. The signature-based assignment is
inconsistent between halves. No permutation we could find transfers.

We note that this test is permissive: estimating a correspondence from the whole
unlabelled target stream is offline source-free adaptation, not test-time adaptation
(Section~\ref{sec:method-tiers}), so it has an advantage no deployable method would have.
Failing under relaxed assumptions implies failing under strict ones.

The conclusion is stronger than the premise it tests. Per-sensor features do not fail
because they are mislabelled across buildings. They fail because a sensor in one
apartment has no counterpart in the other: rooms differ in shape, furniture sits
elsewhere, and routines differ. Relabelling cannot recover information that was never
shared.

\subsection{Three mechanisms}
\label{sec:res-mech-intro}

A natural objection to what follows is that these methods were designed for images, so
their failure on one-dimensional data is to be expected. The results do not support that
reading, for three reasons.

First, the picture is not uniform. DIGA improves on the deployable checkpoint with a
bootstrap interval excluding zero, and AdaBN, the simplest normalisation-based
adaptation there is, gains $41.5$ MCC points in one occupancy direction
(Table~\ref{tab:oe-rev}). Methods built for images are evidently capable of helping on
sensor data; whether they do turns out to depend on the source checkpoint and not on
the modality.

Second, PETAL's failure has nothing to do with modality. Setting $\bar\lambda=0$ deletes
the likelihood term from Eq.~\ref{eq:petal} on any data, images included. What the move to
sensor data changes is only which value a target-label search selects, and even that
follows a statable condition, not the modality.

Third, PeTTA's failure is modality-related but not in the way the objection supposes. The
method does not degrade gracefully; it stops entirely and silently, with losses that look
unremarkable throughout. The mechanism is a scalar that reaches exactly $1$ and multiplies
the update rate by zero. Establishing that is the contribution, not the observation that
image methods sometimes do worse elsewhere.

On implementation: each method is implemented from the equations and hyperparameters given
in its source paper, and we report both the published settings and any retuned ones where
they differ materially. For PeTTA we additionally confirmed the behaviour against the
authors' released implementation, in which the divergence statistic is computed as a sum
over the embedding dimension (\texttt{per\_class\_distance}), which is what produces the
saturation reported below. The PETAL result requires no implementation at all: it follows
from Algorithm~2 of the source paper.
\label{sec:res-mech}

\paragraph{PETAL: the coefficient scales the likelihood.} As noted in
Eq.~\ref{eq:petal}, $\bar\lambda$ weights the data term. Table~\ref{tab:lambda} sweeps it
on the deployable checkpoint with a non-degenerate SWAG posterior.

% source: results/exp12_lambda_bar.csv (src/exp12_lambda_bar.py)
\begin{table*}[t]
\centering
\caption{PETAL as a function of $\bar\lambda$ (published parameterisation,
Eq.~\ref{eq:petal}). Source-only is \macf{} $0.4641$, MCC $0.4679$. At $\bar\lambda=0$
the weights still move, toward the SWAG mean, but the movement is independent of
the target stream.}
\label{tab:lambda}
\small
\begin{tabular}{lcccc}
\toprule
$\bar\lambda$ & Drift from source & Predictions changed & \macf & MCC \\
\midrule
$0.0$  & $1.0021$ & $24/3412$ & $0.4676$ & $0.4724$ \\
$0.1$  & $0.8118$ & $26/3412$ & $0.4650$ & $0.4687$ \\
$1.0$  & $0.3580$ & $17/3412$ & $0.4643$ & $0.4679$ \\
$10.0$ & $0.1688$ & $22/3412$ & $0.4637$ & $0.4677$ \\
\bottomrule
\end{tabular}
\end{table*}

The consequence is specific. On a source model that already outperforms what self-training
produces, selecting $\bar\lambda$ by target-domain \macf{}, a criterion in common use,
prefers $\bar\lambda=0$, the setting at which adaptation cannot occur. A reported null for
PETAL obtained this way is a statement about the search procedure, not about the method.

$\bar\lambda$ is not, however, the whole of why PETAL is inert here, and the table says so:
the predictions barely move at any value, including $\bar\lambda=10$. The second reason is
the length of the stream. This target stream is $3412$ samples, which at batch size $64$
is $54$ gradient steps. PETAL's published learning rate is $10^{-4}$ and its teacher is an
exponential moving average at $\pi=0.999$, so after the entire stream the teacher that
produces the predictions has travelled $1-0.999^{54}=5.3\%$ of the way from the source
weights toward the student. These values are calibrated for the corruption benchmarks the
method was developed on, where a single corruption supplies hundreds of batches and the
full protocol supplies tens of thousands. Fifty-four steps is not a regime in which a
method built on slow teacher-student drift can express itself.

Table~\ref{tab:petal-budget} tests that reading by removing the constraint. Freeing the
teacher and raising the learning rate does make PETAL move, and it moves the wrong way at
every step size.

% source: src/exp22_petal_budget.py
\begin{table}[t]
\centering
\caption{PETAL with the teacher freed ($\pi=0.9$) and the learning rate varied,
$\bar\lambda=1$ throughout, S1$\to$S2. Source-only is \macf{} $0.4641$. The stream is $54$
gradient steps in total.}
\label{tab:petal-budget}
\small
\begin{tabular}{@{}lccc@{}}
\toprule
Learning rate & Predictions changed & \macf & MCC \\
\midrule
$10^{-4}$ (published) & $52/3412$   & $0.4613$ & $0.4667$ \\
$10^{-3}$             & $305/3412$  & $0.4149$ & $0.4112$ \\
$10^{-2}$             & $994/3412$  & $0.3523$ & $0.4006$ \\
$5\times10^{-2}$      & $1767/3412$ & $0.2400$ & $0.1694$ \\
\bottomrule
\end{tabular}
\end{table}

This matters for how the null should be read. PETAL is not merely frozen by a budget that
happens to be too small; when given the budget to act it degrades monotonically. The
method's inactivity is therefore the best outcome available to it on this data, which is a
stronger statement than reporting that nothing happened.

The same selection rule distorts DIGA, in the opposite direction. Running its full
published grid (Section~\ref{sec:background}, $64$ configurations) on the occupancy task,
every one of the $64$ configurations \emph{reduces} MCC relative to source-only once
averaged over the three seeds: the median costs $1.0$ point and the least damaging costs
$0.7$. Across individual runs, $83\%$ are negative and the single best run gains
$0.9$ points, so a search reporting its own maximum over a grid and a seed can return a
positive number drawn from a distribution with no positive mean. Selection on target
labels returns exactly that maximum. A label-free rule, lowest mean prediction entropy
on the target stream, selects the configuration costing $0.7$ points, which is what a
deployment would actually obtain. Section~\ref{sec:res-hgrid} repeats the exercise on the
activity task and for all three methods.

\paragraph{PeTTA: the divergence signal saturates.} Fig.~\ref{fig:saturation} shows why.
Eq.~\ref{eq:petta-gamma} sums over the
$64$-dimensional embedding. Across buildings that sum is large, $\exp(-d)$ underflows, and
$\bar\gamma$ pins at unity, which by Eq.~\ref{eq:petta-alpha} sets $\alpha_t$ to zero and
freezes the teacher. Table~\ref{tab:petta} measures it.

% source: results/exp7_petta_noop.csv (src/exp7_petta_noop.py)
\begin{table*}[t]
\centering
\caption{PeTTA divergence sensing on S2. The published reduction is a sum over embedding
dimensions.}
\label{tab:petta}
\small
\begin{tabular}{lcccc}
\toprule
Variant & median $\bar\gamma$ & median $\alpha_t$ & median $d$ & Changed \\
\midrule
As published (sum)        & $1.000000$ & $4.5\times10^{-9}$ & $587$  & $0/3412$ \\
Per-dimension mean        & $0.8263$   & $8.7\times10^{-3}$ & $9.16$ & $10/3412$ \\
Sum $+$ $\alpha$ floor    & $1.000000$ & $1.0\times10^{-3}$ & $586$  & $1/3412$ \\
\bottomrule
\end{tabular}
\end{table*}


\begin{figure}[t]
\centering
\includegraphics[width=0.92\columnwidth]{fig_saturation.pdf}
\caption{PeTTA's gating function. Because the Mahalanobis statistic is summed over $64$
embedding dimensions, the cross-building median of $587$ sits far inside the saturated
region, giving $\bar\gamma=1$ and an update rate of zero.}
\label{fig:saturation}
\end{figure}

Normalising per dimension unfreezes the mechanism, and it still does not improve on the
source model. An $\alpha$ floor alone does not help, because $\lambda_t=\bar\gamma\lambda_0$
is simultaneously pinned at maximum pull-back. We emphasise that saturation is a failure
mode PeTTA's authors themselves document for ablated variants; what is new here is that it
occurs for the full method when the domain gap is a change of building instead of an
image corruption.

\paragraph{DIGA: published hyperparameters do not transfer.} Table~\ref{tab:diga} crosses
the two readings of DIGA's prototype similarity with two hyperparameter settings.

% source: /tmp/diga_sim.py over tta_faithful.DIGAFaithful
\begin{table*}[t]
\centering
\caption{DIGA under two readings of Eq.~8 and two hyperparameter settings. Source-only is
\macf{} $0.4641$, MCC $0.4679$.}
\label{tab:diga}
\small
\begin{tabular}{llccc}
\toprule
Reading of $\langle\cdot,\cdot\rangle$ & Hyperparameters & Changed & \macf & MCC \\
\midrule
Inner product (literal) & published $0.8/0.9$   & $2573/3412$ & $0.2151$ & $0.1840$ \\
Inner product (literal) & retuned $0.95/0.50$   & $3404/3412$ & $0.0075$ & $-0.1674$ \\
Squared distance        & published $0.8/0.9$   & $1324/3412$ & $0.3932$ & $0.3610$ \\
Squared distance        & retuned $0.95/0.50$   & $357/3412$  & $\mathbf{0.4808}$ & $\mathbf{0.4867}$ \\
Cosine                  & published $0.8/0.9$   & $931/3412$  & $0.3232$ & $0.3399$ \\
Cosine                  & retuned $0.95/0.50$   & $263/3412$  & $0.4750$ & $0.4825$ \\
\bottomrule
\end{tabular}
\end{table*}

The literal inner-product reading yields a negative MCC, confirming that
$\langle\cdot,\cdot\rangle$ must denote a dissimilarity. More substantively, DIGA's
published hyperparameters lose $8.8$ \macf{} points against retuned values on this data.
DIGA's own evaluation makes a point of using a single global parameter set across all
experiments, in explicit contrast to competitors tuned per source--target pair. That
property does not survive the move to one-dimensional sensor data, and recovering the loss
required tuning on target labels, which deployment does not provide.

\subsection{How much of DIGA's gain survives implementation choices?}
\label{sec:res-digaimpl}

Two steps of DIGA are specified in a way that admits more than one implementation, and the
difference is not cosmetic.

The first is the DAM. Eq.~\ref{eq:dam} blends batch-normalisation \emph{statistics}, after
which the input is normalised once. An alternative is to normalise twice, once with source
statistics and once with batch statistics, and blend the two normalised \emph{outputs}. These
are different operations, and a reference implementation in circulation takes the second
route while citing the first equation.

The second is the fusion step. Eq.~\ref{eq:diga-fuse} combines \emph{probabilities}. The
same implementation combines a cosine similarity, bounded in $[-1,1]$, with unbounded
classifier logits, so the relative weight of the two terms depends on logit scale rather
than on $\lambda_F$ alone.

Table~\ref{tab:digaimpl} evaluates all four combinations under the protocol of
Table~\ref{tab:tta}.

% source: src/exp19_diga_impl.py -> results/exp19_diga_impl.csv, 3 seeds, 2000 resamples
\begin{table*}[t]
\centering
\caption{DIGA under four readings of its two under-specified steps, streaming protocol
(13-d layout-invariant features, last GroupKFold fold, no smoothing), 3 seeds.
Source-only is $0.4944$. The four readings share identical source checkpoints, so each
interval is on the paired difference and is narrower than the across-seed spread of any
single column ($\sigma\approx0.022$). No reading yields an interval excluding zero.}
\label{tab:digaimpl}
\small
\begin{tabular}{llcc}
\toprule
DAM & Fusion & \macf & $95\%$ CI on $\Delta$ \\
\midrule
blend statistics (Eq.~\ref{eq:dam}) & probabilities (Eq.~\ref{eq:diga-fuse}) & $\mathbf{0.4996}$ & $[-0.0008,\,+0.0110]$ \\
blend outputs & probabilities & $0.4969$ & $[-0.0043,\,+0.0093]$ \\
blend statistics & similarity $+$ logits & $0.4961$ & $[-0.0035,\,+0.0073]$ \\
blend outputs & similarity $+$ logits & $0.4946$ & $[-0.0058,\,+0.0065]$ \\
\bottomrule
\end{tabular}
\end{table*}

The ordering is the one the equations predict: the reading that follows both published
steps is best and the reading that follows neither is indistinguishable from doing
nothing. But the spread across the four, $0.5$ points, is the same size as the effect
being measured, and no reading separates from zero on this protocol. We report DIGA
throughout using the published equations, which is the reading most favourable to the
method. The point is not that one reading is correct. It is that the adaptation gains at
issue in this study are of a size comparable to the ambiguity in the method's own
description, so a reader reimplementing DIGA from its prose may reasonably obtain no
measurable benefit at all.

\subsection{Does any hyperparameter setting change the conclusion?}
\label{sec:res-hgrid}

The results so far compare each method at its published setting and, where the manuscript
differs, at one alternative. That leaves open whether the activity-task numbers describe
the methods or one point in their configuration space. Table~\ref{tab:hgrid} answers it by
searching each method's own exposed hyperparameters: DIGA over
$\lambda_{BN}\times\mathcal{P}_0\times\rho_P\times K$ ($64$ settings), PETAL over
$\bar\lambda\times\pi\times\delta$ ($24$), and PeTTA over $\alpha_0\times\lambda_0$ and
learning rate ($16$). Each setting is run in both transfer directions with three seeds
under the protocol of Table~\ref{tab:tta}, and is scored two ways: by the best target
\macf{} in the grid, which needs target labels and is therefore an upper bound and not
a result, and by the configuration with the lowest mean prediction entropy on the target
stream, which needs none and is what a deployment can reach.

% source: src/exp20_hparam_grid.py -> results/exp20_hparam_grid.csv
\begin{table*}[t]
\centering
\caption{Full hyperparameter grids on the activity task, $\Delta$\macf{} relative to
source-only, averaged over three seeds. \emph{help} counts settings with a positive mean.
\emph{oracle} is the grid maximum, which requires target labels; \emph{entropy} is the
setting chosen by lowest mean target prediction entropy, which does not. \emph{changed} is
the mean number of predictions altered out of $3412$.}
\label{tab:hgrid}
\small
\begin{tabular}{llccccr}
\toprule
 & & & \multicolumn{3}{c}{$\Delta$\macf} & \\
\cmidrule(lr){4-6}
Method & Direction & help & median & oracle & entropy & changed \\
\midrule
\multicolumn{7}{l}{\emph{source-only} \macf{} $0.4944$ ($\text{S}_1\!\to\!\text{S}_2$), $0.3226$ ($\text{S}_2\!\to\!\text{S}_1$)} \\
\addlinespace[2pt]
DIGA  & $\text{S}_1\!\to\!\text{S}_2$ & $11/64$ & $-0.0176$ & $+0.0112$ & $+0.0031$ & $679$ \\
DIGA  & $\text{S}_2\!\to\!\text{S}_1$ & $64/64$ & $+0.0224$ & $+0.0349$ & $+0.0197$ & $578$ \\
PETAL & $\text{S}_1\!\to\!\text{S}_2$ & $17/24$ & $+0.0002$ & $+0.0028$ & $+0.0028$ & $15$ \\
PETAL & $\text{S}_2\!\to\!\text{S}_1$ & $10/24$ & $-0.0001$ & $+0.0009$ & $-0.0009$ & $20$ \\
PeTTA & $\text{S}_1\!\to\!\text{S}_2$ & $9/16$  & $+0.0000$ & $+0.0003$ & $+0.0000$ & $7$ \\
PeTTA & $\text{S}_2\!\to\!\text{S}_1$ & $4/16$  & $+0.0000$ & $+0.0002$ & $-0.0003$ & $5$ \\
\bottomrule
\end{tabular}
\end{table*}

Three things follow.

For PETAL and PeTTA the grids are flat at zero. No setting in PETAL's published space
moves more than $36$ of $3412$ predictions, and none in PeTTA's moves more than $38$; the
largest target gain anywhere in either grid is $0.3$ points, below the seed spread of the
source model itself. The analysis in Section~\ref{sec:res-mech} attributes this to the
update rules and not to tuning, and the grids are consistent with that: a method whose
data term is switched off by $\bar\lambda=0$, or whose adaptation rate is driven to
$\alpha_t\approx0$ by a saturated gate, does not become active at a different $\pi$ or
$\lambda_0$. Notably, the oracle rule selects $\bar\lambda=0$ for PETAL in the
$\text{S}_1\!\to\!\text{S}_2$ direction: the one value at which, by
Eq.~\ref{eq:petal}, the target stream enters no gradient.

For DIGA the grid is wide, and which way it points is decided by the source model rather
than by the hyperparameters. Transferring from the stronger source, $53$ of $64$ settings
lose and the median costs $1.8$ points; transferring from the weaker one, all $64$ gain
and the median adds $2.2$. Placing this beside the occupancy grid of
Section~\ref{sec:res-mech} gives the same ordering across both tasks:

\begin{center}\footnotesize
\begin{tabular}{@{}lcc@{}}
\toprule
Source & Quality & Helped \\
\midrule
occupancy $\text{S}_1$ & MCC $0.6583$ & $0/64$ \\
activity $\text{S}_1$  & \macf{} $0.4944$ & $11/64$ \\
activity $\text{S}_2$  & \macf{} $0.3226$ & $64/64$ \\
\bottomrule
\end{tabular}
\end{center}

The fraction of a published grid that helps tracks how much headroom the direction has.
Entering S2 the cross-building penalty is $0.6$ points (Section~\ref{sec:res-gap}) and
$53$ of $64$ settings lose; entering S1 it is $21.3$ points and all $64$ gain. Read over
whole configuration spaces rather than at single points, normalisation-statistic
adaptation is
not a property of a method that one tunes into working, it is a trade of source-fitted
statistics for target-fitted ones, and whether that trade pays depends on which of the two
is better.

Third, the two selection rules disagree, and the direction of the disagreement matters for
how such results should be read. Label-free selection recovers $28\%$ of DIGA's oracle gain
in the forward direction and $56\%$ in the reverse; for PETAL and PeTTA it returns a
negative number in the reverse direction where the oracle returns a positive one. An
entropy criterion prefers confident predictions, and a method can raise confidence while
lowering accuracy, so the gap is expected, not anomalous. Its size is the point: the
headroom a grid search reports is between a quarter and a half larger than what a system
without target labels can select, which is the condition TTA is defined under.

\subsection{Occupancy estimation}
\label{sec:res-oe}

Table~\ref{tab:oe} runs the identical model and methods under three normalisation
protocols.

% source: results/exp13_occ_scaler.csv (src/exp13_occ_scaler.py), 5 seeds
\begin{table*}[t]
\centering
\caption{Occupancy estimation, MCC on the target test set ($n=150$, $34$ minority), under
three normalisation protocols. Always predicting ``not occupied'' gives MCC $0.000$.}
\label{tab:oe}
\begin{tabular}{lccc}
\toprule
Method & Fitted on source & Fitted on target train & Fitted on target train $+$ test \\
\midrule
Source-only & $\mathbf{0.6583}\pm0.013$ & $\mathbf{0.6951}\pm0.013$ & $\mathbf{0.6966}\pm0.017$ \\
AdaBN       & $0.6412$ & $0.6232$ & $0.6323$ \\
DIGA        & $0.6401$ & $0.6693$ & $0.6769$ \\
PETAL       & $0.6583$ & $0.6951$ & $0.6966$ \\
PeTTA       & $0.6583$ & $0.6951$ & $0.6966$ \\
\bottomrule
\end{tabular}
\end{table*}

Three results follow from Table~\ref{tab:oe}, all of them about the forward direction;
the reverse is taken up below. No method improves on the source model under any of the
three normalisation protocols, and DIGA is below the baseline in all three columns. The choice of normalisation is worth $+3.8$ MCC
points, but it is worth that to \emph{every} method equally, including the unadapted
baseline. It is normalisation, not adaptation, and a protocol that fits the scaler on
target statistics will credit it to whichever method happens to be under test. Finally,
PETAL and PeTTA reproduce
source-only predictions to four decimal places on this task as well, which rules out a
HAR-specific explanation for the two freezes reported in Section~\ref{sec:res-mech}.

Running the transfer in the reverse direction, training on the target office and
evaluating on the source office, changes the picture completely, and in a way that is
the most direct evidence in this paper for its central claim.

% source: results/exp17_occ_reverse.csv (src/exp17_occ_reverse.py), 5 seeds
\begin{table*}[t]
\centering
\caption{Occupancy in both transfer directions, MCC, source-fitted normalisation, 5 seeds.
The only thing that differs between the two columns is which office supplies the source
model. Always predicting ``not occupied'' gives MCC $0.000$ in both.}
\label{tab:oe-rev}
\small
\begin{tabular}{lcc}
\toprule
Method & source $\to$ target & target $\to$ source \\
\midrule
Source-only & $\mathbf{0.6583}$ & $-0.0775$ \\
AdaBN       & $0.6412$ & $\mathbf{0.3374}$ \\
DIGA        & $0.6401$ & $0.1934$ \\
PETAL       & $0.6583$ & $-0.0775$ \\
PeTTA       & $0.6583$ & $-0.0775$ \\
\bottomrule
\end{tabular}
\end{table*}

In the reverse direction the source checkpoint is not merely weak but worse than chance:
trained on the target office, it reaches MCC $-0.0775$ on the source office, meaning its
predictions are anti-correlated with the labels. Against that starting point AdaBN
improves MCC by $41.5$ points and DIGA by $27.1$. In the forward direction, starting from
a checkpoint at $0.6583$, the same two methods lose $1.7$ and $1.8$ points respectively.

The contrast is as clean as this dataset allows: identical methods, identical protocol,
identical features, two offices, and the only variable is whether the source model was any
good. When it is, adaptation has nothing to contribute and the normalisation-based methods
actively harm it; when it is broken, the same methods recover a large part of the deficit.
This mirrors AdaBN's behaviour on HAR, where it costs $17$ \macf{} points against a sound
checkpoint, and it supports reading adaptation gain as a measure of source-model
deficiency and not of adaptation quality. A method evaluated only against a weak
source model will appear to work.

PETAL and PeTTA reproduce the source predictions exactly in this direction too. Across two
tasks and two directions each, neither method has changed a prediction.

% ============================================================================
\section{Discussion}
\label{sec:discussion}

The headline result of this study is not about any one method. It is that the quantity
these methods are asked to repair, the cross-building performance drop, is routinely
measured in a way that substantially overstates it, and that most of what remains once it is
measured consistently is attributable to feature design and not to anything TTA can
reach.

Measured against the right ceiling, the cross-building penalty is $0.6$ points entering
one apartment and $21.3$ entering the other (Section~\ref{sec:res-gap}), against the $83$
points implied by comparing a leakage-inflated source figure with an honest target one.
In the forward direction a model built from thirteen layout-invariant statistics performs
on the second apartment as well as a model trained on that apartment's own days. The
in-domain ceiling then becomes the binding constraint: one cannot improve target
performance much beyond it, which is why our attempts to push the source model further
(richer invariant features, class weighting, wider networks) gained little. Only causal
smoothing helped, and modestly.

The asymmetry is what we would most want a reader to take away, because it decides what
adaptation can do. Where the penalty is $0.6$ points there is no headroom, and across
DIGA's full published grid $53$ of $64$ settings lose ground. Where it is $21.3$ points
all $64$ gain. Occupancy estimation reproduces the same split on a second task and a
different model family. A reported adaptation gain is therefore first of all a statement
about the gap the method was given, and only afterwards about the method.

This also reframes what per-sensor features cost, and what in-domain validation is worth.
In-domain those features are valuable: the best source-validation score in
Table~\ref{tab:features} belongs to a set that retains them. Across buildings they are
worse than useless, scoring below chance alone in both directions. The consequence for
model selection is real but narrower than the extremes of those tables suggest. Pooled
over all $70$ configurations in this study, source validation and target performance are
positively correlated ($\rho=+0.63$, $p<0.001$); among competitive configurations the
relationship weakens to $\rho=+0.27$ ($p=0.11$), and the best-validating configuration is
not the best-transferring one in either direction. In-domain validation reliably separates
working configurations from broken ones. Between reasonable alternatives, which is the
choice a practitioner actually faces, it carries little information about which will
survive a change of building, and a held-out building is what settles it.

On the methods themselves, the three failures are instructive precisely because they are
different from one another. PETAL's is a reading failure with a sharp edge: a coefficient
that looks like a regulariser weight scales the likelihood, and the hyperparameter search
most practitioners would run, selecting on target performance, actively prefers the
value that disables adaptation. PeTTA's is dimensional: a divergence statistic that behaves
sensibly when the domain gap is a corruption of the same image saturates when the gap is a
different building, and because the mechanism fails silently the loss curves look
unremarkable throughout. DIGA's is about generality: it works, but the single global
parameter set it advertises does not, and the gap between its published and retuned
settings is larger than the gain adaptation produces.

A practical consequence runs through all three. Hyperparameter selection on held-out target
labels, which is standard in this literature, is not merely optimistic. It can select degenerate configurations, as the PETAL sweep shows, and it can
make a method appear to transfer when what transferred was the tuning. A label-free
criterion should be reported alongside any oracle-tuned number.

The occupancy results in both directions (Table~\ref{tab:oe-rev}) are the cleanest
demonstration of this available here, because they hold everything constant but the
quality of the source model. Starting from a checkpoint at MCC $0.658$, adaptation costs
roughly two points; starting from one at $-0.078$, the same methods gain twenty-seven to
forty-one. The methods did not change. What changed is how much room there was.

\paragraph{What generalises, and what does not.} The findings above differ in how far
they travel, and it is worth separating them.

Two do not depend on this dataset at all. PETAL's $\bar\lambda$ multiplies the
cross-entropy in its published objective (Eq.~\ref{eq:petal}), so $\bar\lambda=0$ removes
the data term on any dataset, in any modality; The dataset-specific part is
narrower than it may appear: selection on target labels prefers $\bar\lambda=0$ whenever
the source model already outperforms what self-training produces, which is a condition
that can be checked, not a coincidence of this data. PeTTA's gating function
likewise saturates whenever the summed Mahalanobis statistic is large, which is a property
of summing over embedding dimensions rather than of these apartments, though whether
the statistic reaches that regime does depend on how far apart the domains are, and the
authors' own traces on image corruptions sit at $0.8$--$1.0$ rather than at unity.

The evaluation result is a property of a procedure, not of a dataset: stride-$1$
windowing with a random validation split places near-duplicates across the split whenever
windows overlap, which is common in sensor-based HAR. That overlapping windows bias
cross-validation in activity recognition is established \cite{hammerla2015}; what we add
is the magnitude on this pipeline, $46$ points of source \macf{}, and the separation of
the split from the window definition, which contributes under two.

The asymmetry of the cross-building penalty is different in kind. It is measured here on
two apartments, so the specific values of $0.6$ and $21.3$ points belong to this pair and
nothing follows about their magnitude elsewhere. What we would expect to generalise is the
procedure and its consequence: a cross-building penalty should be measured against a model
trained on the target building at matched sample size, not against source validation, and
until it is measured that way an adaptation gain cannot be separated from the headroom it
was handed. Establishing how often the penalty is negligible requires datasets with more
than two environments.

\paragraph{Limitations.} The study rests on one HAR dataset and two apartments from it.
Running both transfer directions rules out a direction-specific fluke but says nothing
about other buildings, other sensor modalities, or richer activity sets; external validity
beyond these two apartments is untested, and the dataset contains no further homes. The
occupancy data is private, so that section is not independently
reproducible from the released code. Our best configuration improves on the simpler one by
$3.3$ \macf{} points against a between-fold standard deviation of $0.067$, and we present
it as a modest gain and not a result. And the \texttt{cleaning} class, with $113$ target
samples, is not learned by any method we tried; whether that is a sample-size limit or a
representational one, we cannot say from these data.

% ============================================================================
\section{Conclusion}
\label{sec:conclusion}

We evaluated three image-domain test-time adaptation methods on cross-building human
activity recognition and occupancy estimation, and found that the evaluation protocol
dominated the comparison. Stride-one windowing with a random validation split inflates
source-domain \macf{} from roughly $0.53$ to $0.998$; grouping the split by recording day
removes the effect, and the window definition is nearly irrelevant to it. Making features
causal, which the online premise of TTA requires, improves target transfer by $7.3$ points
rather than costing anything.

Measured against a model trained on the target building at matched sample size, the
cross-building penalty is $0.6$ points entering one apartment and $21.3$ entering the
other. One of our two directions therefore contains essentially no domain shift, and what
the conventional source-minus-target figure counts as an $8.8$-point gap there is mostly
the difference between the two apartments as prediction problems. Most of the apparent gap
that does exist is attributable to features encoding sensor identity, which score below
stratified-random guessing on the target in both directions.

The methods follow that asymmetry. Across DIGA's full published grid, $53$ of $64$
settings lose ground in the direction with no headroom and all $64$ gain in the direction
with $21$ points of it; its best reported figure, $1.7$ \macf{} points, requires
hyperparameters retuned on target labels, and a label-free criterion recovers $0.3$.
PETAL and PeTTA do not move at all, for reasons traceable to their own specifications.
PETAL's published learning rate and teacher update rate assume streams of thousands of
batches, and ours is $54$; forcing it to move by raising either makes it monotonically
worse. PeTTA's divergence signal sums over the embedding dimension and saturates when the
domains are buildings. Both reproduce on occupancy estimation, in both of its transfer
directions, where the two directions differ only in the quality of the source checkpoint:
against a sound one no method improves, while against one worse than chance the same
methods gain $27$ to $41$ MCC points. A further apparent gain proves attributable to a
target-fitted normalisation worth $3.8$ MCC points to every method, the unadapted one
included.

We do not conclude that test-time adaptation is unsuitable for smart buildings. We conclude
that it was being asked to repair a problem largely created by evaluation and feature
design, and that fixing those leaves little for adaptation to contribute. Future work
should look at making building-specific information transferable instead of discarding
it, aligning sensors across sites by functional role using unlabelled co-activation
structure, since the points separating our layout-invariant features from the
best in-domain configuration are exactly the information currently being thrown away.

% ============================================================================
\section*{Reproducibility}

Code, the raw sensor data and every result file are available at
\url{https://github.com/mariemkbaier/TTA-Test-Time-Adaptation}. The repository reproduces the
full analysis from the raw CSVs in one pass, each table above names the script and result
file that generated it. The occupancy dataset is private: it
was collected at the Grenoble Institute of Technology and is not ours to redistribute,
so that section alone is reported but not independently reproducible. The three adaptation
methods are implemented in \texttt{src/tta.py} (as commonly applied) and
\texttt{src/tta\_faithful.py} (at the hyperparameters and equations each paper publishes);
the corrected PETAL of Section~\ref{sec:res-mech} is \texttt{src/petal\_fixed.py}.

\section*{Acknowledgements}

This research was supported by the Natural Sciences and Engineering Research Council of
Canada (NSERC).

% ============================================================================
\appendix
\section{A larger layout-invariant feature set}
\label{app:inv}

We also evaluated a $60$-dimensional invariant set adding a sorted activation profile (the
per-sensor window means sorted in descending order, which is invariant to any permutation
of sensor indices), run-length and recency structure, and the window statistics recomputed
at $W\in\{5,10,40\}$. It does not help:

\begin{table*}[t]
\centering
\caption{Richer layout-invariant features, S1$\to$S2, GroupKFold, 3 seeds.}
\begin{tabular}{lcc}
\toprule
Feature set & $d$ & Target \macf \\
\midrule
Global $+$ temporal                      & 13 & $\mathbf{0.4350}\pm0.061$ \\
$+$ sorted activation profile            & 28 & $0.4245\pm0.039$ \\
$+$ run-length structure                 & 33 & $0.4227\pm0.027$ \\
$+$ multi-scale windows                  & 60 & $0.4188\pm0.028$ \\
\bottomrule
\end{tabular}
\end{table*}

Permutation invariance in sensor index is evidently not sufficient: S1 has $70$ sensors and
S2 has $75$, placed in differently shaped rooms, so ``how many sensors are busy and how
concentrated the activity is'' is itself building-specific. Inverse-frequency class
weighting also failed, reducing \macf{} to $0.3852$ and \emph{lowering} balanced accuracy
from $0.4818$ to $0.4280$, most plausibly because it pushes the model toward
\texttt{cleaning}, which is not learnable from $113$ target samples, at the cost of the
four classes that do work.

% ============================================================================
\begin{thebibliography}{99}

\bibitem{hammerla2015} N. Y. Hammerla, T. Pl\"otz, Let's (not) stick together: pairwise
similarity biases cross-validation in activity recognition, in: Proceedings of the 2015
ACM International Joint Conference on Pervasive and Ubiquitous Computing, 2015,
pp. 1041--1051.

\bibitem{manfren2021} M. Manfren, M. Sibilla, L. Tronchin, Energy modelling and analytics
in the built environment---a review of their role for energy transitions in the
construction sector, Energies 14 (3) (2021) 679.

\bibitem{amangeldy2025} B. Amangeldy, T. Imankulov, N. Tasmurzayev, G. Dikhanbayeva,
Y. Nurakhov, A review of artificial intelligence and deep learning approaches for resource
management in smart buildings, Buildings 15 (15) (2025) 2631.

\bibitem{dridi2022} J. Dridi, M. Amayri, N. Bouguila, Transfer learning for estimating
occupancy and recognizing activities in smart buildings, Building and Environment 217
(2022) 109057.

\bibitem{amayri2019} M. Amayri, S. Ploix, N. Bouguila, F. Wurtz, Estimating occupancy using
interactive learning with a sensor environment: real-time experiments, IEEE Access 7 (2019)
53932--53944.

\bibitem{dridi_cluster} J. Dridi, M. Amayri, N. Bouguila, Unsupervised clustering-based
domain adaptation for estimating occupancy and recognizing activities in smart buildings,
Journal of Building Engineering 85 (2024) 108741.

\bibitem{shahbazian2023} R. Shahbazian, I. Trubitsyna, Human sensing by using radio
frequency signals: a survey on occupancy and activity detection, IEEE Access 11 (2023)
40878--40904.

\bibitem{hong2017} T. Hong, D. Yan, S. D'Oca, C. Chen, Ten questions concerning occupant
behavior in buildings: the big picture, Building and Environment 114 (2017) 518--530.

\bibitem{liang2024} J. Liang, R. He, T. Tan, A comprehensive survey on test-time adaptation
under distribution shifts, International Journal of Computer Vision 133 (1) (2025) 31--64.

\bibitem{diga} W. Wang, Z. Zhong, W. Wang, X. Chen, C. Ling, B. Wang, N. Sebe, Dynamically
instance-guided adaptation: a backward-free approach for test-time domain adaptive semantic
segmentation, in: CVPR, 2023, pp. 24090--24099.

\bibitem{petal} D. Brahma, P. Rai, A probabilistic framework for lifelong test-time
adaptation, in: CVPR, 2023, pp. 3582--3591.

\bibitem{petta} T.-H. Hoang, D. M. Vo, M. N. Do, Persistent test-time adaptation in
recurring testing scenarios, in: NeurIPS, Vol. 37, 2024.

\bibitem{tent} D. Wang, E. Shelhamer, S. Liu, B. Olshausen, T. Darrell, Tent: fully
test-time adaptation by entropy minimization, in: ICLR, 2021.

\bibitem{cotta} Q. Wang, O. Fink, L. Van Gool, D. Dai, Continual test-time domain
adaptation, in: CVPR, 2022, pp. 7201--7211.

\bibitem{shot} J. Liang, D. Hu, J. Feng, Do we really need to access the source data?
Source hypothesis transfer for unsupervised domain adaptation, in: ICML, Vol. 119 of PMLR,
2020, pp. 6028--6039.

\bibitem{abderrahim2025} O. Abderrahim, J. Dridi, M. Amayri, N. Bouguila, Occupancy
estimation and activity recognition in smart buildings using open set domain adaptation,
Journal of Building Engineering 112 (2025) 113784.

\bibitem{abderrahim2026} O. Abderrahim, J. Dridi, M. Amayri, N. Bouguila, Scalable activity
recognition in smart buildings via generalized domain adaptation of IoT sensor data, Energy
and Buildings 351 (2026) 116692.

\bibitem{abderrahim_osda} O. Abderrahim, J. Dridi, M. Amayri, N. Bouguila, Open set domain
adaptation and universal domain adaptation for enhancing activity recognition in smart
buildings, in: Proceedings of the 19th IBPSA Conference, Brisbane, Australia, 2025.
\url{https://doi.org/10.26868/25222708.2025.1331}.

\bibitem{wang2023} S. Wang, J. Wang, H. Xi, B. Zhang, L. Zhang, H. Wei, Optimization-free
test-time adaptation for cross-person activity recognition, Proceedings of the ACM on
Interactive, Mobile, Wearable and Ubiquitous Technologies 7 (4) (2023) 183.

\bibitem{tapia2004} E. M. Tapia, S. S. Intille, K. Larson, Activity recognition in the
home using simple and ubiquitous sensors, in: Pervasive Computing, Springer, Berlin,
Heidelberg, 2004, pp. 158--175.

\bibitem{chicco2020} D. Chicco, G. Jurman, The advantages of the Matthews correlation
coefficient (MCC) over F1 score and accuracy in binary classification evaluation, BMC
Genomics 21 (1) (2020) 6.

\end{thebibliography}

\end{document}
